\chapter{Technical reference}

\section{DTMF $\to$ keyboard mapping}

\begin{table}[h]
\centering
\begin{tabular}{ccccp{5cm}}
\toprule
\textbf{DTMF} & \textbf{Frequencies (Hz)} & \textbf{Key} & \textbf{evdev code} & \textbf{Action in Hugo} \\
\midrule
1 & 697 + 1209 & Q & KEY\_Q & free / action 1 \\
2 & 697 + 1336 & $\uparrow$ & KEY\_UP & jump or up \\
3 & 697 + 1477 & E & KEY\_E & free / action 2 \\
4 & 770 + 1209 & $\leftarrow$ & KEY\_LEFT & left \\
5 & 770 + 1336 & Space & KEY\_SPACE & main action \\
6 & 770 + 1477 & $\rightarrow$ & KEY\_RIGHT & right \\
7 & 852 + 1209 & Z & KEY\_Z & free / action 3 \\
8 & 852 + 1336 & $\downarrow$ & KEY\_DOWN & down / crouch \\
9 & 852 + 1477 & X & KEY\_X & free / action 4 \\
* & 941 + 1209 & Left Shift & KEY\_LEFTSHIFT & modifier \\
0 & 941 + 1336 & Esc & KEY\_ESC & pause / menu \\
\# & 941 + 1477 & Enter & KEY\_ENTER & confirm \\
\bottomrule
\end{tabular}
\caption{Default DTMF $\to$ keyboard mapping.}
\end{table}

\section{Key system files}

\begin{table}[h]
\centering
\small
\begin{tabular}{p{0.45\textwidth} p{0.45\textwidth}}
\toprule
\textbf{File} & \textbf{Purpose} \\
\midrule
\texttt{/etc/asterisk/sip.conf} & Defines SIP extensions 101 and 102 \\
\texttt{/etc/asterisk/extensions.conf} & Dialplan: logic for 9000 \\
\texttt{/etc/asterisk/manager.conf} & Enables AMI with the \texttt{hugo-bridge} user \\
\texttt{/etc/systemd/system/hugo-bridge.service} & Bridge service unit \\
\texttt{/etc/udev/rules.d/99-hugo-uinput.rules} & Permissions on \texttt{/dev/uinput} \\
\texttt{/opt/hugo-phone/bridge/hugo-bridge} & Compiled bridge binary \\
\texttt{/opt/hugo-phone/bridge/config.toml} & Bridge configuration (hot-reload) \\
\texttt{/opt/hugo-phone/games/hugo/} & Game files \\
\texttt{\textasciitilde/.config/dosbox/dosbox-staging.conf} & Emulator configuration \\
\bottomrule
\end{tabular}
\end{table}

\section{Useful services and commands}

\subsection{Managing Asterisk}

\begin{shellcode}
sudo systemctl status asterisk
sudo systemctl restart asterisk
sudo asterisk -rvvv               # interactive console
sudo asterisk -rx "reload"        # reload configs without restart
\end{shellcode}

Commands inside the Asterisk CLI:

\begin{plaincode}
sip show peers              ; extension states
sip show registry           ; registered clients
sip set debug on/off        ; detailed traces
dialplan show hugo-game     ; show loaded dialplan
manager show settings       ; AMI state
manager show users          ; AMI users
manager show connected      ; active AMI sessions
core show channels          ; active calls
\end{plaincode}

\subsection{Managing the bridge}

\begin{shellcode}
sudo systemctl status hugo-bridge
sudo systemctl restart hugo-bridge
sudo journalctl -u hugo-bridge -f      # live logs
sudo journalctl -u hugo-bridge -n 100  # last 100 lines
\end{shellcode}

\subsection{Rebuilding the bridge}

\begin{shellcode}
cd /opt/hugo-phone/bridge
sudo -u beto ~/.cargo/bin/cargo build --release
sudo cp target/release/hugo-bridge ./hugo-bridge
sudo systemctl restart hugo-bridge
\end{shellcode}

\subsection{Inspecting the virtual keyboard}

\begin{shellcode}
ls /dev/input/by-id/                   # look for hugo-phone-...
sudo evtest                            # select the device
                                       # watch events when pressing DTMF
\end{shellcode}

\subsection{HT801 health check}

\begin{shellcode}
ping 192.168.1.11                      # IP reachability
nmap -p 5060 192.168.1.11              # SIP port open
\end{shellcode}

\section{Structure of the \texttt{config.toml} file}

\begin{tomlcode}
[ami]
host = "127.0.0.1"
port = 5038
username = "hugo-bridge"
secret = "REAL-PASSWORD"

[bridge]
press_duration_ms = 80
event_name = "HugoKey"
key_field = "Key"

[keymap]
"2" = "Up"
"4" = "Left"
"6" = "Right"
"8" = "Down"
"5" = "Space"
"0" = "Esc"
"1" = "Q"
"3" = "E"
"7" = "Z"
"9" = "X"
"*" = "LeftShift"
"#" = "Enter"
\end{tomlcode}

\section{Supported key names}

The \texttt{keyboard.rs} parser recognizes:

\begin{itemize}
  \item \textbf{Directional:} \texttt{Up}, \texttt{Down}, \texttt{Left}, \texttt{Right}
  \item \textbf{Special:} \texttt{Space}, \texttt{Enter}, \texttt{Esc} (or \texttt{Escape}), \texttt{Tab}, \texttt{Backspace}
  \item \textbf{Modifiers:} \texttt{LeftShift}, \texttt{RightShift}, \texttt{LeftCtrl}, \texttt{RightCtrl}, \texttt{LeftAlt}, \texttt{RightAlt}
  \item \textbf{Letters:} \texttt{A} through \texttt{Z}
  \item \textbf{Function keys:} \texttt{F1} through \texttt{F12}
\end{itemize}

If an unknown name is assigned, the bridge logs a warning but keeps running.

\section{AMI protocol anatomy}

\subsection{Initial banner}

\begin{plaincode}
Asterisk Call Manager/9.0.0
\end{plaincode}

\subsection{Login}

\begin{plaincode}
Action: Login
Username: hugo-bridge
Secret: REAL-PASSWORD
Events: user,call
(blank line)
\end{plaincode}

\subsection{Successful login response}

\begin{plaincode}
Response: Success
Message: Authentication accepted
(blank line)
\end{plaincode}

\subsection{UserEvent produced by the dialplan}

\begin{plaincode}
Event: UserEvent
Privilege: user,all
UserEvent: HugoKey
Channel: PJSIP/101-00000003
Caller: 101
Key: 2
Uniqueid: 1715641234.3
(blank line)
\end{plaincode}

\section{Glossary}

\begin{description}
  \item[AMI] \emph{Asterisk Manager Interface}. Plain-text TCP protocol for controlling Asterisk from external clients.
  \item[ATA] \emph{Analog Telephone Adapter}. Device that converts an analog line to SIP.
  \item[DOSBox] Emulator of the MS-DOS operating system and the IBM PC compatible hardware. DOSBox-Staging is a modernized fork.
  \item[DTMF] \emph{Dual-Tone Multi-Frequency}. Signaling system using pairs of tones, used by telephone keypads.
  \item[evdev] Linux kernel subsystem that exposes input devices to \emph{user space}.
  \item[FXO] \emph{Foreign eXchange Office}. Port that connects to a switchboard as if it were a telephone.
  \item[FXS] \emph{Foreign eXchange Subscriber}. Port that supplies line as if it were the phone company.
  \item[KMSDRM] \emph{Kernel Mode Setting / Direct Rendering Manager}. Linux kernel layer that allows drawing directly to the screen without an X server.
  \item[PBX] \emph{Private Branch Exchange}. Private telephone switchboard.
  \item[RFC 2833] Standard for carrying DTMF tones inside RTP packets, out of band relative to the main audio.
  \item[RTP] \emph{Real-time Transport Protocol}. Protocol for transmitting audio and video in real time, over UDP.
  \item[SIP] \emph{Session Initiation Protocol}. Signaling protocol for initiating, maintaining and terminating voice/video sessions.
  \item[uinput] Linux kernel subsystem that lets \emph{user space} create virtual input devices.
\end{description}

\section{External references}

\begin{itemize}
  \item Official Asterisk documentation: \texttt{wiki.asterisk.org}
  \item DOSBox-Staging: \texttt{dosbox-staging.github.io}
  \item Rust \texttt{uinput} crate: \texttt{docs.rs/uinput}
  \item Tokio: \texttt{tokio.rs}
  \item Grandstream HT801 manual: \texttt{grandstream.com/products}
  \item Raspberry Pi documentation: \texttt{raspberrypi.com/documentation}
  \item Hugo on MyAbandonware: \texttt{myabandonware.com/game/hugo-eya}
\end{itemize}

\section{Credits}

Hugo, the character and the franchise, are the property of \textbf{ITE Media ApS} (originally \emph{Interactive Television Entertainment}), Denmark. This project is a non-profit technical tribute with no official affiliation with the company.

The complete system, its code and its documentation were developed as a personal project by \textbf{Beto} in collaboration with Claude (Anthropic) during May 2026.

\vspace{1cm}

\begin{center}
\itshape\small
``Hugo, watch out for that rock!''\\
--- Generations of Chilean children in the nineties
\end{center}
